% P8 future work + closing
\section{Future Work: Building the Hybrid}
\label{sec:p8-future}

The obvious next step is to build and measure the architecture rather than
argue it. Our plan, in the order the seats de-risk each other:

\begin{enumerate}[leftmargin=1.5em, itemsep=1pt, topsep=2pt]
\item \textbf{Sensor first.} Attach a VLM feature extractor to a
document-bearing subset of cases and measure the marginal AUC of
sensor-derived features \emph{inside the existing tree}---the cleanest test,
since the metric is the incumbent's own. This includes adversarial
evaluation: injected instructions inside documents must degrade at worst that
document's features \citep{greshake2023injection}.
\item \textbf{Scribe second.} Evaluate SAR-narrative drafts against FinCEN's
element-coverage structure (who/what/when/where/why/how)
\citep{fincen2003sarnarrative} with analyst-grader rubrics, and
adverse-action drafts for fidelity to the tree's actual attributions
\citep{cfr1002_9, cfpb2022circular}---measuring time-to-file and revision
distance, with a hard human sign-off gate.
\item \textbf{Triage third.} Replace the internal $85\times$ estimate with
measured per-alert token costs and analyst-minutes-saved on a live queue,
including the ordering quality of agent-ranked versus score-ranked queues.
\item \textbf{Cold-start last.} Simulate typology emergence by holding out a
generator cluster as an ``unseen scheme'' and measuring few-shot LLM detection
during the label vacuum, against the retrained tree's catch-up curve.
\end{enumerate}

Each experiment scores one seat with the seat's own metric---capability
measured, label ignored. If the program thesis holds anywhere outside RL
post-training, it should hold here: the pipeline that results is not ``AI
replacing the model'' but a tree that scores, an LLM that sees and writes, and
an analyst who decides.
